\IfFileExists{stacks-project.cls}{%
\documentclass{stacks-project}
}{%
\documentclass{amsart}
}

% For dealing with references we use the comment environment
\usepackage{verbatim}
\newenvironment{reference}{\comment}{\endcomment}
%\newenvironment{reference}{}{}
\newenvironment{slogan}{\comment}{\endcomment}
\newenvironment{history}{\comment}{\endcomment}

% For commutative diagrams we use Xy-pic
\usepackage[all]{xy}

% We use 2cell for 2-commutative diagrams.
\xyoption{2cell}
\UseAllTwocells

% We use multicol for the list of chapters between chapters
\usepackage{multicol}

% This is generally recommended for better output
\usepackage{lmodern}
\usepackage[T1]{fontenc}

% For cross-file-references
\usepackage{xr-hyper}

% Package for hypertext links:
\usepackage{hyperref}

% For any local file, say "hello.tex" you want to link to please
% use \externaldocument[hello-]{hello}
\externaldocument[introduction-]{introduction}
\externaldocument[conventions-]{conventions}
\externaldocument[sets-]{sets}
\externaldocument[categories-]{categories}
\externaldocument[topology-]{topology}
\externaldocument[sheaves-]{sheaves}
\externaldocument[sites-]{sites}
\externaldocument[stacks-]{stacks}
\externaldocument[fields-]{fields}
\externaldocument[algebra-]{algebra}
\externaldocument[brauer-]{brauer}
\externaldocument[homology-]{homology}
\externaldocument[derived-]{derived}
\externaldocument[simplicial-]{simplicial}
\externaldocument[more-algebra-]{more-algebra}
\externaldocument[smoothing-]{smoothing}
\externaldocument[modules-]{modules}
\externaldocument[sites-modules-]{sites-modules}
\externaldocument[injectives-]{injectives}
\externaldocument[cohomology-]{cohomology}
\externaldocument[sites-cohomology-]{sites-cohomology}
\externaldocument[dga-]{dga}
\externaldocument[dpa-]{dpa}
\externaldocument[sdga-]{sdga}
\externaldocument[hypercovering-]{hypercovering}
\externaldocument[schemes-]{schemes}
\externaldocument[constructions-]{constructions}
\externaldocument[properties-]{properties}
\externaldocument[morphisms-]{morphisms}
\externaldocument[coherent-]{coherent}
\externaldocument[divisors-]{divisors}
\externaldocument[limits-]{limits}
\externaldocument[varieties-]{varieties}
\externaldocument[topologies-]{topologies}
\externaldocument[descent-]{descent}
\externaldocument[perfect-]{perfect}
\externaldocument[more-morphisms-]{more-morphisms}
\externaldocument[flat-]{flat}
\externaldocument[groupoids-]{groupoids}
\externaldocument[more-groupoids-]{more-groupoids}
\externaldocument[etale-]{etale}
\externaldocument[chow-]{chow}
\externaldocument[intersection-]{intersection}
\externaldocument[pic-]{pic}
\externaldocument[weil-]{weil}
\externaldocument[adequate-]{adequate}
\externaldocument[dualizing-]{dualizing}
\externaldocument[duality-]{duality}
\externaldocument[discriminant-]{discriminant}
\externaldocument[derham-]{derham}
\externaldocument[local-cohomology-]{local-cohomology}
\externaldocument[algebraization-]{algebraization}
\externaldocument[curves-]{curves}
\externaldocument[resolve-]{resolve}
\externaldocument[models-]{models}
\externaldocument[functors-]{functors}
\externaldocument[equiv-]{equiv}
\externaldocument[pione-]{pione}
\externaldocument[etale-cohomology-]{etale-cohomology}
\externaldocument[proetale-]{proetale}
\externaldocument[relative-cycles-]{relative-cycles}
\externaldocument[more-etale-]{more-etale}
\externaldocument[trace-]{trace}
\externaldocument[crystalline-]{crystalline}
\externaldocument[spaces-]{spaces}
\externaldocument[spaces-properties-]{spaces-properties}
\externaldocument[spaces-morphisms-]{spaces-morphisms}
\externaldocument[decent-spaces-]{decent-spaces}
\externaldocument[spaces-cohomology-]{spaces-cohomology}
\externaldocument[spaces-limits-]{spaces-limits}
\externaldocument[spaces-divisors-]{spaces-divisors}
\externaldocument[spaces-over-fields-]{spaces-over-fields}
\externaldocument[spaces-topologies-]{spaces-topologies}
\externaldocument[spaces-descent-]{spaces-descent}
\externaldocument[spaces-perfect-]{spaces-perfect}
\externaldocument[spaces-more-morphisms-]{spaces-more-morphisms}
\externaldocument[spaces-flat-]{spaces-flat}
\externaldocument[spaces-groupoids-]{spaces-groupoids}
\externaldocument[spaces-more-groupoids-]{spaces-more-groupoids}
\externaldocument[bootstrap-]{bootstrap}
\externaldocument[spaces-pushouts-]{spaces-pushouts}
\externaldocument[spaces-chow-]{spaces-chow}
\externaldocument[groupoids-quotients-]{groupoids-quotients}
\externaldocument[spaces-more-cohomology-]{spaces-more-cohomology}
\externaldocument[spaces-simplicial-]{spaces-simplicial}
\externaldocument[spaces-duality-]{spaces-duality}
\externaldocument[formal-spaces-]{formal-spaces}
\externaldocument[restricted-]{restricted}
\externaldocument[spaces-resolve-]{spaces-resolve}
\externaldocument[formal-defos-]{formal-defos}
\externaldocument[defos-]{defos}
\externaldocument[cotangent-]{cotangent}
\externaldocument[examples-defos-]{examples-defos}
\externaldocument[algebraic-]{algebraic}
\externaldocument[examples-stacks-]{examples-stacks}
\externaldocument[stacks-sheaves-]{stacks-sheaves}
\externaldocument[criteria-]{criteria}
\externaldocument[artin-]{artin}
\externaldocument[quot-]{quot}
\externaldocument[stacks-properties-]{stacks-properties}
\externaldocument[stacks-morphisms-]{stacks-morphisms}
\externaldocument[stacks-limits-]{stacks-limits}
\externaldocument[stacks-cohomology-]{stacks-cohomology}
\externaldocument[stacks-perfect-]{stacks-perfect}
\externaldocument[stacks-introduction-]{stacks-introduction}
\externaldocument[stacks-more-morphisms-]{stacks-more-morphisms}
\externaldocument[stacks-geometry-]{stacks-geometry}
\externaldocument[moduli-]{moduli}
\externaldocument[moduli-curves-]{moduli-curves}
\externaldocument[examples-]{examples}
\externaldocument[exercises-]{exercises}
\externaldocument[guide-]{guide}
\externaldocument[desirables-]{desirables}
\externaldocument[coding-]{coding}
\externaldocument[obsolete-]{obsolete}
\externaldocument[fdl-]{fdl}
\externaldocument[index-]{index}

% Theorem environments.
%
\theoremstyle{plain}
\newtheorem{theorem}[subsection]{Theorem}
\newtheorem{proposition}[subsection]{Proposition}
\newtheorem{lemma}[subsection]{Lemma}

\theoremstyle{definition}
\newtheorem{definition}[subsection]{Definition}
\newtheorem{example}[subsection]{Example}
\newtheorem{exercise}[subsection]{Exercise}
\newtheorem{situation}[subsection]{Situation}

\theoremstyle{remark}
\newtheorem{remark}[subsection]{Remark}
\newtheorem{remarks}[subsection]{Remarks}

\numberwithin{equation}{subsection}

% Macros
%
\def\lim{\mathop{\mathrm{lim}}\nolimits}
\def\colim{\mathop{\mathrm{colim}}\nolimits}
\def\Spec{\mathop{\mathrm{Spec}}}
\def\Hom{\mathop{\mathrm{Hom}}\nolimits}
\def\Ext{\mathop{\mathrm{Ext}}\nolimits}
\def\SheafHom{\mathop{\mathcal{H}\!\mathit{om}}\nolimits}
\def\SheafExt{\mathop{\mathcal{E}\!\mathit{xt}}\nolimits}
\def\Sch{\mathit{Sch}}
\def\Mor{\mathop{\mathrm{Mor}}\nolimits}
\def\Ob{\mathop{\mathrm{Ob}}\nolimits}
\def\Sh{\mathop{\mathit{Sh}}\nolimits}
\def\NL{\mathop{N\!L}\nolimits}
\def\CH{\mathop{\mathrm{CH}}\nolimits}
\def\proetale{{pro\text{-}\acute{e}tale}}
\def\etale{{\acute{e}tale}}
\def\QCoh{\mathit{QCoh}}
\def\Ker{\mathop{\mathrm{Ker}}}
\def\Im{\mathop{\mathrm{Im}}}
\def\Coker{\mathop{\mathrm{Coker}}}
\def\Coim{\mathop{\mathrm{Coim}}}

% Boxtimes
%
\DeclareMathSymbol{\boxtimes}{\mathbin}{AMSa}{"02}

%
% Macros for moduli stacks/spaces
%
\def\QCohstack{\mathcal{QC}\!\mathit{oh}}
\def\Cohstack{\mathcal{C}\!\mathit{oh}}
\def\Spacesstack{\mathcal{S}\!\mathit{paces}}
\def\Quotfunctor{\mathrm{Quot}}
\def\Hilbfunctor{\mathrm{Hilb}}
\def\Curvesstack{\mathcal{C}\!\mathit{urves}}
\def\Polarizedstack{\mathcal{P}\!\mathit{olarized}}
\def\Complexesstack{\mathcal{C}\!\mathit{omplexes}}
% \Pic is the operator that assigns to X its picard group, usage \Pic(X)
% \Picardstack_{X/B} denotes the Picard stack of X over B
% \Picardfunctor_{X/B} denotes the Picard functor of X over B
\def\Pic{\mathop{\mathrm{Pic}}\nolimits}
\def\Picardstack{\mathcal{P}\!\mathit{ic}}
\def\Picardfunctor{\mathrm{Pic}}
\def\Deformationcategory{\mathcal{D}\!\mathit{ef}}
\usepackage{cleveref}
\crefname{originallemma}{original lemma}{original lemmas}
\Crefname{originallemma}{Original lemma}{Original lemmas}
\makeatletter
\newlabel{lemma-jump-loci}{{0BDI}{source}{}{sourcepointer-1}{}}
\newlabel{lemma-jump-loci@cref}{{[originallemma][0][]0BDI}{[1][0][]source}{}{}{}}
\makeatother
\crefname{originallemma}{original lemma}{original lemmas}
\Crefname{originallemma}{Original lemma}{Original lemmas}
\makeatletter
\newlabel{lemma-perfect-affine}{{08EB}{source}{}{sourcepointer-2}{}}
\newlabel{lemma-perfect-affine@cref}{{[originallemma][0][]08EB}{[1][0][]source}{}{}{}}
\makeatother
\crefname{originallemma}{original lemma}{original lemmas}
\Crefname{originallemma}{Original lemma}{Original lemmas}
\makeatletter
\newlabel{more-algebra-lemma-lift-perfect-from-residue-field}{{0BCD}{source}{}{sourcepointer-3}{}}
\newlabel{more-algebra-lemma-lift-perfect-from-residue-field@cref}{{[originallemma][0][]0BCD}{[1][0][]source}{}{}{}}
\makeatother
\crefname{originallemma}{original lemma}{original lemmas}
\Crefname{originallemma}{Original lemma}{Original lemmas}
\makeatletter
\newlabel{more-algebra-lemma-coker-injective-free}{{0EWY}{source}{}{sourcepointer-4}{}}
\newlabel{more-algebra-lemma-coker-injective-free@cref}{{[originallemma][0][]0EWY}{[1][0][]source}{}{}{}}
\makeatother
\crefname{originallemma}{original lemma}{original lemmas}
\Crefname{originallemma}{Original lemma}{Original lemmas}
\makeatletter
\newlabel{lemma-affine-compare-bounded}{{06Z0}{source}{}{sourcepointer-5}{}}
\newlabel{lemma-affine-compare-bounded@cref}{{[originallemma][0][]06Z0}{[1][0][]source}{}{}{}}
\makeatother
\title{The locus where a cohomology module of a perfect complex is locally free of a given rank is locally closed}
\author{The Stacks Project Authors}
\date{stacks-project:perfect.tex:local-snapshot-20261012; GNU FDL 1.2 or later}
\begin{document}
\maketitle
\setcounter{subsection}{1}
\noindent\textit{Editorial note.} All delivered mathematical text is the author's own, copied from the pinned Stacks Project chapter file perfect.tex; no statement or proof step was rewritten, reordered or removed. The counted span is the author environment and its own proof environment. References inside the span are handed to the assembler as navigation stubs that carry the canonical tag of the cited tag only; whether each of them is load-bearing is recorded on the stub itself.\par\medskip
\noindent\textit{Allowed standard prerequisites (editorial).}\begin{itemize}
\item The level sets of the dimension function of the fibre cohomology are locally constructible and the function is upper semicontinuous. invoked at proof lines 8081-8095 and again at proof lines 8096-8100. Produces the open where the dimension is at most the given rank and shows the set where it is smaller is open.
\item A perfect complex is locally represented by a bounded complex of finite free modules with ranks prescribed by the fibres. invoked at proof lines 8100-8110. Supplies the local model used to construct the ideal cutting out the locus.
\item A map of finite free modules whose cokernel is injective and free gives a locally closed condition described by the vanishing of minors. invoked at proof lines 8110-8134. Produces the finite type ideal defining the closed subscheme and the criterion in terms of the rank of the cohomology module.
\end{itemize}\medskip
\ifdefined\hypertarget\hypertarget{sourcepointer-1}{}\fi
\noindent\textit{Source reference.} Original lemma 0BDI: sources/stacks-perfect/pinned-original-source.tex, source lines 7907--7907. Invoked three times, at proof lines 8081-8084, 8089-8092 and 8093-8095, for the properties of the dimension function; every conclusion drawn from them is inside the counted span.\par\medskip
\ifdefined\hypertarget\hypertarget{sourcepointer-2}{}\fi
\noindent\textit{Source reference.} Original lemma 08EB: sources/stacks-perfect/pinned-original-source.tex, source lines 2636--2636. Cited at proof lines 8100-8108 for the local representation used to write down the matrix; the ideal defined from that matrix is inside the counted span.\par\medskip
\ifdefined\hypertarget\hypertarget{sourcepointer-3}{}\fi
\noindent\textit{Source reference.} Original lemma 0BCD: sources/stacks-chow/tags.txt, source lines 1--1. Used at proof lines 8103-8110 to choose the ranks in the local model; the construction built on it is inside the counted span.\par\medskip
\ifdefined\hypertarget\hypertarget{sourcepointer-4}{}\fi
\noindent\textit{Source reference.} Original lemma 0EWY: sources/stacks-chow/tags.txt, source lines 1--1. Invoked at proof lines 8113-8125 to translate the condition on the cokernel into the vanishing of minors, which is how the ideal is defined; the analysis of the resulting subscheme is inside the counted span.\par\medskip
\ifdefined\hypertarget\hypertarget{sourcepointer-5}{}\fi
\noindent\textit{Source reference.} Original lemma 06Z0: sources/stacks-perfect/pinned-original-source.tex, source lines 454--454. A navigation pointer in the reduction to the affine local model; the reduction itself is inside the counted span.\par\medskip
\begin{lemma}
\label{lemma-locally-closed-where-H0-locally-free}
Let $X$ be a scheme. Let $E \in D(\mathcal{O}_X)$ be perfect
of tor-amplitude in $[a, b]$ for some $a, b \in \mathbf{Z}$.
Let $r \geq 0$.
Then there exists a locally closed subscheme $j : Z \to X$
characterized by the following
\begin{enumerate}
\item $H^a(Lj^*E)$ is a locally free $\mathcal{O}_Z$-module of rank $r$, and
\item a morphism $f : Y \to X$ factors through $Z$
if and only if for all morphisms $g : Y' \to Y$ the
$\mathcal{O}_{Y'}$-module $H^a(L(f \circ g)^*E)$ is locally free
of rank $r$.
\end{enumerate}
Moreover, $j : Z \to X$ is of finite presentation and we have
\begin{enumerate}
\item[(3)] if $f : Y \to X$ factors as $Y \xrightarrow{g} Z \to X$, then
$H^a(Lf^*E) = g^*H^a(Lj^*E)$,
\item[(4)] if $\beta_a(x) \leq r$ for all $x \in X$, then
$j$ is a closed immersion and given $f : Y \to X$ the following
are equivalent
\begin{enumerate}
\item $f : Y \to X$ factors through $Z$,
\item $H^a(Lf^*E)$ is a locally free $\mathcal{O}_Y$-module of rank $r$,
\end{enumerate}
and if $r = 1$ these are also equivalent to
\begin{enumerate}
\item[(c)] $\mathcal{O}_Y \to \SheafHom_{\mathcal{O}_Y}(H^a(Lf^*E), H^a(Lf^*E))$
is injective.
\end{enumerate}
\end{enumerate}
\end{lemma}

\begin{proof}
First, let $U \subset X$ be the locally constructible
open subscheme where the function
$\beta_a$ of Lemma \ref{lemma-jump-loci} has values $\leq r$.
Let $f : Y \to X$ be as in (2). Then for any $y \in Y$
we have $\beta_a(Lf^*E) = r$ hence $y$ maps into $U$ by
Lemma \ref{lemma-jump-loci}. Hence $f$ as in (2) factors through $U$.
Thus we may replace $X$ by $U$ and assume that
$\beta_a(x) \in \{0, 1, \ldots, r\}$ for all $x \in X$.
We will show that in this case
there is a closed subscheme $Z \subset X$ cut out by a finite type
quasi-coherent ideal characterized
by the equivalence of (4) (a), (b) and (4)(c) if $r = 1$
and that (3) holds.
This will finish the proof because it will a fortiori show
that morphisms as in (2) factor through $Z$.

\medskip\noindent
If $x \in X$ and $\beta_a(x) < r$, then there is an
open neighbourhood of $x$ where $\beta_a < r$
(Lemma \ref{lemma-jump-loci}). In this way we
see that set theoretically at least $Z$ is a closed subset.

\medskip\noindent
To get a scheme theoretic structure, consider a point $x \in X$
with $\beta_a(x) = r$. Set $\beta = \beta_{a + 1}(x)$.
By More on Algebra, Lemma
\ref{more-algebra-lemma-lift-perfect-from-residue-field}
there exists an affine open neighbourhood $U$ of $x$
such that $K|_U$ is represented by a complex
$$
\ldots \to 0 \to \mathcal{O}_U^{\oplus r}
\xrightarrow{(f_{ij})} \mathcal{O}_U^{\oplus \beta} \to
\ldots \to
\mathcal{O}_U^{\oplus \beta_{b - 1}(x)} \to
\mathcal{O}_U^{\oplus \beta_b(x)} \to 0
\to \ldots
$$
(This also uses earlier results to turn the problem into algebra, for example
Lemmas \ref{lemma-affine-compare-bounded} and
\ref{lemma-perfect-affine}.) Now, if $g : Y \to U$ is any morphism
of schemes such that $g^\sharp(f_{ij})$ is nonzero for some pair $i, j$, then
$H^a(Lg^*E)$ is not a locally free $\mathcal{O}_Y$-module of rank $r$.
See More on Algebra, Lemma \ref{more-algebra-lemma-coker-injective-free}.
Trivially $H^a(Lg^*E)$ is a locally free $\mathcal{O}_Y$-module if
$g^\sharp(f_{ij}) = 0$ for all $i, j$. Thus we see that over $U$ the
closed subscheme cut out by all $f_{ij}$ satisfies
(3) and we have the equivalence of (4)(a) and (b).
The characterization of $Z$ shows that the locally constructed patches
glue (details omitted). Finally, if $r = 1$ then
(4)(c) is equivalent to (4)(b) because in this case locally
$H^a(Lg^*E) \subset \mathcal{O}_Y$ is the annihilator of the ideal generated
by the elements $g^\sharp(f_{ij})$.
\end{proof}

\end{document}
